\documentclass[]{../../../../tex/classes/styledArticle}
\usepackage{../../../../tex/packages/hebrewSupport}
\usepackage{../../../../tex/packages/mathShortcuts}
\usepackage{../../../../tex/packages/theoremsSupport}
\usepackage{../../../../tex/packages/enfancysections}

%! ~~~ Document ~~~

\author{שחר פרץ}
\title{\textit{לינארית 2א 20}}
\date{22 ביוני 2025}
\begin{document}
	\maketitle
	תזכורת: $A \in M_n(\C)$ נקראית סימטרית אממ $A = A^T$ והרמיטית אם $A = A^*$, ונורמלית אממ $AA^* = A^*A$. 
	
	\theo{תהי $T \co V \to V$ ט''ל, ו־$V$ ממ''פ מעל $\F \in \{\R, \C\}$, ויהי $B$ בסיס א''נ של $V$. אזי אם $A = [T]_B$: 
	\[ [T^*]_B = A = ([T]_B)^* \]}  
	\begin{proof}
		נזכר ש־: 
		\[ [T]_B = \pms{\vert &  & \vert \\ [Te_1]_B & \cdots & [Te_n]_B \\ \vert &  & \vert} \]
		נסמן $B = \{e_i\}_{i = 1}^{n}$ בסיס. נבחין ש־: 
		\[ Te_j = \sumni a_{ij}e_i, \ a_{ij} = \mut{Te_j}{e_i} \]
		נסמן ב־$C$ את המטריצה המייצגת $[T^*]_B$: 
		\[ c_{ij} = \mut{T^*e_j}{e_i} \]
		ונחשב: 
		\[ c_ij = \mut{T^*e_j}{e_i} = \mut{e_j}{Te_i} = \ol{\mut{Te_i}{e_j}} = a_{ij} \]
		
	\end{proof}
	
	\textbf{מסקנה: }אם $A$ נורמלית אז $T_A$ נורמלית מעל $\F^n$ אם הסטנדרטית. בפרט מתקיים עליה המשפט הספקטרלי. גם אם $A$ ממשית, הע''ע עלולים להמצא מעל $\C$ (אלא אם היא צמודה לעצמה, ואז הם מעל $\R$). 
	
	משהו על אינטרפולציות: 
	\theo{יהיו $x_1 \dots x_n, y_1 \dots y_n \in \R$. נניח $\forall i, j \in [n] \co i \neq j \implies x_i \neq x_j$. אז $\exists! p \in \R_{\le n - 1}[x] \co \forall i \in [n] \co p(x_i) = y_i$ עד לכדי חברות (באופן שקול: נניח $p$ מתוקן) }
	(הערה מהידע הכללי שלי: זהו פולינום לגראנג' והוא בונה אינטרפולציה די נחמדה). \begin{proof}
		ידוע שהפולינום מהצורה $p(x)  = \sum_{k = 0}^{n - 1} a_k x^k  = (1, x, x^2, \dots x^{n - 1}) (a_0 \dots a_{n - 1})^T$. למעשה, נקבל את מטריצת ונדרמונד: 
		\[ \pms{1 & x_1 & x_1^2 & \cdots & x_1^n \\ 1 & x_2 & x_2^2 & \cdots & x_2^{n - 1} \\ \vdots & \vdots & \vdots & \ddots & \vdots \\ 1 & x_n & x_n^2 & \cdots & x_n^{n - 1}}\pms{a_0 \\ a_1 \\ \vdots \\ a_{n - 1}} = \pms{y_1 \\ y_2 \\ \vdots \\ y_n} \]
		וידוע שהדטרמיננטה של ונדרמונד היא $\prod_{i < j} (x_i - x_j)$ וזה איכשהו אמור לגמור את ההוכחה. 
		
		אם $x_i = y_i$ בפולינום לעיל, אז $\forall a \in \C \co f(\bar a) = \ol{f(a)} \implies f \in \R[x]$. הוכחה: נניח בשלילה, אז $\exists \ag \in \C\setminus \R \co f(\ag) = 0$ כך ש־$f(\bar \ag) = 0$. אזי $0 \neq f(\bar \ag) = \ol{f(\ag)}  = 0$ וזו סתירה. 
	\end{proof}
	
	\theo{תהי $A \in M_n(\F)$ נורמלית, אז קיים פולינום $f(x) \in \R[x]$ כך ש־$A^* = f(A)$. }
	\textit{הערה: }באופן כללי \textbf{לא} נכון שאם $A, B$ מתחלפות אז $\exists f(x) \in \F[x] \co f(A) = B$. 
	\begin{proof}
		נוכיח מעל $\C$ ובפרט נכונות ל־$\R$. מהמשפט הספקטרלי קיימת $P$ הפיכה (מעבר לבסיס אורתונורמלי) כך ש־$P\op A P = \diag(\lg_1 \dots \lg_n)$ הוא הבסיס המלכסן. מהיות המלכסן א''נ נקבל: 
		\[ P\op A^* P = \diag(\lg_1 \dots \lg_n)^* = \diag(\bar\lg_1 \dots \bar \lg_n) \]
		ולכן מפולינום האינטרפולציה עבור הקבוצות $\{\lg_1 \dots \lg_n\}$ ו־$\{\bar \lg_1 \dots \bar \lg_n\}$ קיים ויחיד עד לכדי חברות $f(x) \in \R[x]$ כך ש־$f(\lg_i) = \bar \lg_i$. נותר להראות ש־$A^* = f(A)$, זאת כי: 
		\[ P\op A^* P = f(\diag(\lg_1 \dots \lg_n)) = \diag(\bar \lg_1 \dots \bar \lg_n) \implies f(P\op A P) = P\op f(A) P \]
		נכפול ב־$P$ מצד אחד וב־$P\op$ מהצד השני ונקבל את הדרוש
	\end{proof}
	
	
	\ndoc
\end{document}
